\chapter{Kết luận}
\label{chap:conclusion}

\section{Kết luận}
\label{sec:conclusion}

Luận văn này nghiên cứu bài toán phân loại lối sống thực khuẩn thể từ dữ liệu metagenomic, một thách thức quan trọng trong sinh tin học hiện đại. Các phương pháp hiện có gặp khó khăn khi xử lý contig ngắn do phụ thuộc vào công cụ dự đoán gene như Prodigal~\cite{hyatt2010prodigal}, vốn không hoạt động hiệu quả trên các đoạn DNA ngắn thiếu ORF hoàn chỉnh~\cite{trimble2012short}. Để giải quyết vấn đề này, luận văn đề xuất PhaBERT-CNN, một kiến trúc học sâu lai kết hợp mô hình nền tảng genome DNABERT-2~\cite{zhou2023dnabert} với mô đun Modified TextCNN gồm nhánh CNN đa tỷ lệ và nhánh attention pooling. Phương pháp hoạt động trực tiếp trên chuỗi DNA thô thông qua BPE tokenization, không phụ thuộc vào công cụ dự đoán gene, và được huấn luyện theo chiến lược hai giai đoạn với discriminative learning rates để tối ưu hóa cả backbone và các thành phần task-specific.

Phương pháp được đánh giá toàn diện trên bốn nhóm độ dài contig (100--400 bp, 400--800 bp, 800--1200 bp, 1200--1800 bp) thông qua stratified 5-fold cross-validation trên tập dữ liệu gồm 2,241 hệ gen thực khuẩn thể hoàn chỉnh. Như trình bày trong Bảng~\ref{tab:main_results}, PhaBERT-CNN đạt hiệu suất tốt nhất trên ba nhóm contig ngắn và trung bình (A, B, C) với accuracy lần lượt là 81.59\%, 87.91\%, và 90.01\%. So với các phương pháp tiên tiến, PhaBERT-CNN cải thiện 0.91--2.67\% so với PhaTYP~\cite{shang2023phatyp}, 2.17--3.24\% so với DeePhage~\cite{wu2021deephage}, và 0.91--5.98\% so với ProkBERT~\cite{ligeti2024prokbert} trên các nhóm này. Trên contig dài (Nhóm D), PhaBERT-CNN duy trì kết quả cạnh tranh với 90.69\% accuracy, chỉ kém PhaTYP 0.34\% nhưng vẫn vượt trội DeePhage 1.60\% và ProkBERT 5.96\%. Đặc biệt, PhaBERT-CNN đạt specificity cao nhất ở Nhóm D (90.95\%), vượt trội PhaTYP (89.75\%), cho thấy khả năng phát hiện phage ôn hòa tốt hơn trên contig dài.

Kết quả nghiên cứu chứng minh tính hiệu quả của việc kết hợp mô hình nền tảng genome tiền huấn luyện với kiến trúc chuyên biệt task-specific cho bài toán phân loại lối sống thực khuẩn thể. PhaBERT-CNN không chỉ cải thiện độ chính xác phân loại trên contig ngắn và trung bình mà còn giải quyết các hạn chế quan trọng của phương pháp hiện có như phụ thuộc gene prediction và khả năng tổng quát hóa hạn chế. Sự kết hợp giữa biểu diễn ngữ cảnh toàn cục từ DNABERT-2 và trích xuất đặc trưng cục bộ chuyên biệt từ CNN đa tỷ lệ tạo ra lợi thế rõ rệt so với từng thành phần riêng lẻ, mở ra hướng tiếp cận đầy hứa hẹn cho việc áp dụng genomic foundation models vào các bài toán sinh tin học downstream khác như dự đoán vi khuẩn chủ, phát hiện prophage, và phân loại genome vi sinh vật.

\section{Thảo luận về kết quả chính}
\label{sec:discussion}

\subsection{Phân tích kết quả}
\label{sec:result_analysis}

PhaBERT-CNN đạt hiệu suất vượt trội trên ba nhóm contig ngắn và trung bình (A, B, C) với accuracy từ 81.59\% đến 90.01\%, cao hơn đáng kể so với các phương pháp baseline. Hiệu suất này xuất phát từ ba yếu tố bổ sung cho nhau. Thứ nhất, PhaBERT-CNN hoạt động trực tiếp trên chuỗi DNA thô thông qua BPE tokenization, không phụ thuộc vào công cụ dự đoán gene như Prodigal mà PhaTYP yêu cầu. Trên contig ngắn, Prodigal thường không phát hiện được ORF hoàn chỉnh~\cite{trimble2012short}, khiến PhaTYP mất đi nguồn thông tin chính; trong khi đó PhaBERT-CNN vẫn khai thác được tín hiệu từ chuỗi nucleotide thô. Thứ hai, DNABERT-2 đã được tiền huấn luyện trên corpus genome đa loài quy mô lớn~\cite{zhou2023dnabert}, cho phép mô hình tổng quát hóa tốt hơn DeePhage vốn được huấn luyện từ đầu trên tập dữ liệu hạn chế. Thứ ba, kiến trúc nhánh song song kép—CNN đa tỷ lệ với kernel sizes 3, 5, 7 kết hợp attention pooling—cho phép mô hình đồng thời nắm bắt các motif cục bộ ở nhiều tỷ lệ không gian và tổng hợp ngữ cảnh toàn cục có trọng số, tạo ra biểu diễn phong phú hơn so với chỉ dùng DNABERT-2 đơn thuần.

Kết quả so sánh giữa DNABERT-2 baseline và PhaBERT-CNN (Hình~\ref{fig:dnabert2_vs_phabertcnn}) cho thấy các thành phần task-specific đóng góp đáng kể vào hiệu suất tổng thể. Cải thiện về accuracy tăng dần theo độ dài contig—từ 0.07\% ở Nhóm A (81.52\% lên 81.59\%) đến 3.43\% ở Nhóm D (87.26\% lên 90.69\%)—phản ánh rằng khi chuỗi dài hơn, các motif chức năng trở nên rõ ràng hơn và CNN có thể khai thác chúng hiệu quả hơn. Đặc biệt, cải thiện specificity ở Nhóm A đạt 6.90\% (từ 73.25\% lên 80.15\%), cho thấy attention pooling giúp mô hình tập trung vào các vùng thông tin nhất của chuỗi ngắn, giảm false positive trong phát hiện phage ôn hòa. Kết quả này xác nhận giả thuyết rằng việc kết hợp pre-trained embeddings với kiến trúc chuyên biệt task-specific mang lại lợi ích rõ rệt, đặc biệt khi các thành phần này được thiết kế để bổ sung cho nhau: transformer encoder cung cấp biểu diễn ngữ cảnh toàn cục, CNN trích xuất các mẫu cục bộ đa tỷ lệ, và attention pooling tổng hợp thông tin có trọng số.

Trên Nhóm D (1200--1800 bp), PhaTYP vượt PhaBERT-CNN với accuracy 91.02\% so với 90.69\%, một khoảng cách nhỏ chỉ 0.34\%. Hiệu suất vượt trội của PhaTYP trên contig dài xuất phát từ khả năng tận dụng protein features từ các ORF hoàn chỉnh có mặt trong chuỗi dài, cung cấp tín hiệu bổ sung về chức năng sinh học mà PhaBERT-CNN chưa khai thác. Tuy nhiên, PhaBERT-CNN đạt specificity cao nhất ở Nhóm D (90.95\%), vượt trội PhaTYP (89.75\%) và tất cả các phương pháp khác, cho thấy khả năng phát hiện phage ôn hòa chính xác hơn trên contig dài. Điều này gợi ý rằng việc kết hợp thêm thông tin protein vào PhaBERT-CNN theo kiến trúc multi-modal—xử lý song song cả nucleotide sequence và protein features—có thể cải thiện hiệu suất trên contig dài trong khi vẫn duy trì ưu thế trên contig ngắn, tạo ra một phương pháp robust trên toàn bộ phổ độ dài contig.

So sánh với các phương pháp còn lại cho thấy rõ ưu thế của việc sử dụng genomic foundation models. DeePhage thể hiện hiệu suất trung bình với accuracy từ 78.07\% đến 89.09\%, cải thiện đều đặn theo độ dài contig nhưng luôn thấp hơn các phương pháp dựa trên transformer. Hạn chế này xuất phát từ việc DeePhage sử dụng CNN trên mã hóa one-hot, không tận dụng được kiến thức từ pre-training và chỉ học từ tập dữ liệu huấn luyện hạn chế. ProkBERT có xu hướng đặc biệt: accuracy đạt đỉnh ở Nhóm B (85.35\%) rồi giảm xuống khoảng 84\% ở Nhóm C và D, gợi ý rằng tokenization LCA chồng lấn của ProkBERT gặp khó khăn trong việc nắm bắt các phụ thuộc tầm xa trong chuỗi dài hơn. Hiện tượng stagnation này phản ánh hạn chế cơ bản của k-mer tokenization chồng lấn, vốn có thể gây information leakage và không mở rộng hiệu quả cho chuỗi dài~\cite{zhou2023dnabert}. Ngược lại, PhaBERT-CNN duy trì hiệu suất ổn định và cải thiện đều đặn trên tất cả nhóm độ dài, cho thấy BPE tokenization và kiến trúc lai của DNABERT-2 giải quyết tốt hơn các thách thức này.

\subsection{Đóng góp của nghiên cứu}
\label{sec:contributions}

Về mặt phương pháp, luận văn đề xuất PhaBERT-CNN, một kiến trúc học sâu lai mới kết hợp mô hình nền tảng genome DNABERT-2 với mô đun Modified TextCNN chuyên biệt cho bài toán phân loại lối sống thực khuẩn thể. Kiến trúc nhánh song song kép bao gồm nhánh CNN đa tỷ lệ với ba kernel sizes (3, 5, 7) để trích xuất các motif cục bộ ở nhiều tỷ lệ không gian, và nhánh attention pooling để tổng hợp biểu diễn toàn cục có trọng số từ các token embeddings của DNABERT-2. Chiến lược huấn luyện hai giai đoạn với discriminative learning rates—giai đoạn 1 đóng băng backbone để huấn luyện classification head, giai đoạn 2 fine-tune toàn bộ mô hình với learning rate thấp hơn cho backbone—cho phép tối ưu hóa hiệu quả cả pre-trained representations và task-specific components. Phương pháp này khác biệt với các tiếp cận trước đây bằng cách hoạt động trực tiếp trên chuỗi DNA thô thông qua BPE tokenization, loại bỏ phụ thuộc vào gene prediction tools và cho phép xử lý hiệu quả các contig ngắn thiếu ORF hoàn chỉnh.

Về mặt kết quả thực nghiệm, luận văn cung cấp đánh giá toàn diện trên bốn nhóm độ dài contig từ 100 đến 1800 bp thông qua stratified 5-fold cross-validation, cho thấy PhaBERT-CNN cải thiện state-of-the-art trên contig ngắn và trung bình với mức tăng accuracy 0.91--5.98\% so với các phương pháp tiên tiến. Nghiên cứu chứng minh rằng việc kết hợp pre-trained genomic embeddings với task-specific architecture mang lại lợi ích rõ rệt: so sánh trực tiếp với DNABERT-2 baseline cho thấy các thành phần CNN và attention pooling đóng góp cải thiện accuracy từ 0.07\% đến 3.43\% tùy theo độ dài contig, với cải thiện specificity đặc biệt đáng kể (6.90\%) trên contig ngắn. Kết quả này xác nhận giả thuyết rằng genomic foundation models cung cấp nền tảng mạnh mẽ cho các tác vụ downstream, nhưng vẫn cần các thành phần chuyên biệt để đạt hiệu suất tối ưu. Phân tích chi tiết về hiệu suất trên từng nhóm độ dài cũng cung cấp insights quan trọng về điểm mạnh và hạn chế của từng phương pháp, giúp định hướng nghiên cứu tương lai.

Về mặt ứng dụng thực tiễn, PhaBERT-CNN giải quyết bài toán quan trọng trong phân tích metagenomic: phân loại lối sống thực khuẩn thể từ các contig phân mảnh thu được từ metagenomic assemblies. Khả năng hoạt động trực tiếp trên chuỗi DNA thô mà không phụ thuộc gene prediction tools làm cho phương pháp đặc biệt phù hợp với các kịch bản thực tế, nơi contig thường ngắn và thiếu protein-coding regions hoàn chỉnh. Hiệu suất cân bằng giữa sensitivity và specificity trên tất cả nhóm độ dài cho thấy PhaBERT-CNN có thể được tích hợp vào các pipeline phân tích metagenomic tự động để hỗ trợ nghiên cứu về tương tác phage-bacteria, động lực quần thể vi sinh vật, và ứng dụng phage therapy. Hơn nữa, kiến trúc lai kết hợp foundation model với task-specific components đặt ra một paradigm mới có thể được áp dụng cho các bài toán sinh tin học downstream khác như dự đoán vi khuẩn chủ, phát hiện prophage, phân loại chức năng gene, và annotation genome vi sinh vật, mở ra nhiều hướng nghiên cứu tiềm năng trong lĩnh vực computational genomics.

\subsection{Hạn chế và hướng phát triển tương lai}
\label{sec:limitations_future}

Một hạn chế quan trọng của PhaBERT-CNN là việc sử dụng random undersampling để xử lý class imbalance trong tập dữ liệu huấn luyện. Phương pháp này, mặc dù đơn giản và hiệu quả về mặt tính toán, có thể gây mất thông tin từ lớp đa số (virulent phages) khi loại bỏ ngẫu nhiên các mẫu để cân bằng với lớp thiểu số (temperate phages). Các chiến lược thay thế như class weighting—gán trọng số cao hơn cho lớp thiểu số trong loss function—hoặc focal loss—tập trung vào các mẫu khó phân loại—có thể giữ lại toàn bộ thông tin huấn luyện trong khi vẫn đảm bảo mô hình không bị bias về lớp đa số. Oversampling tổng hợp sử dụng data augmentation techniques như reverse complement, random mutation với tỷ lệ thấp, hoặc mixup trên embedding space cũng có thể được khảo sát để tăng cường tập dữ liệu lớp thiểu số mà không làm mất thông tin lớp đa số. Nghiên cứu tương lai nên đánh giá hệ thống các chiến lược này để xác định phương pháp tối ưu cho bài toán phân loại lối sống thực khuẩn thể.

Hạn chế thứ hai liên quan đến hiệu suất trên contig dài, nơi PhaTYP vượt trội PhaBERT-CNN nhờ tận dụng protein features từ các ORF hoàn chỉnh. Điều này gợi ý rằng thông tin protein cung cấp tín hiệu bổ sung quan trọng cho phân loại lối sống, đặc biệt trên chuỗi dài chứa các protein-coding regions hoàn chỉnh. Hướng phát triển tự nhiên là mở rộng PhaBERT-CNN theo kiến trúc multi-modal, xử lý song song cả nucleotide sequence thông qua DNABERT-2 và protein features thông qua protein language models như ESM-2 hoặc ProtBERT. Kiến trúc multi-modal này có thể sử dụng cross-attention mechanism để tích hợp thông tin từ hai modalities, hoặc late fusion để kết hợp predictions từ hai nhánh độc lập. Quan trọng hơn, kiến trúc cần được thiết kế để hoạt động robust khi protein features không có sẵn (trên contig ngắn), có thể thông qua conditional computation hoặc learned gating mechanisms. Cách tiếp cận này có tiềm năng cải thiện hiệu suất trên contig dài trong khi vẫn duy trì ưu thế trên contig ngắn, tạo ra một phương pháp truly universal cho toàn bộ phổ độ dài contig.

Hạn chế thứ ba là phạm vi đánh giá hiện tại chỉ tập trung vào simulated contigs được tạo từ complete genomes thông qua sliding window. Mặc dù phương pháp này cho phép đánh giá có kiểm soát trên các nhóm độ dài cụ thể, nó có thể không phản ánh đầy đủ các thách thức trong real metagenomic data như sequencing errors, chimeric contigs, horizontal gene transfer, và prophage regions. Nghiên cứu tương lai nên mở rộng đánh giá sang các tập dữ liệu metagenomic thực tế từ các môi trường đa dạng như human gut, marine, soil, và wastewater để kiểm tra khả năng tổng quát hóa của PhaBERT-CNN. Việc so sánh hiệu suất trên simulated vs. real data cũng sẽ cung cấp insights quan trọng về gap giữa controlled evaluation và practical application, giúp định hướng cải tiến phương pháp cho các kịch bản thực tế.

Hướng phát triển quan trọng khác là tăng cường interpretability của PhaBERT-CNN để hiểu rõ hơn cơ chế phân loại của mô hình. Phân tích attention weights từ attention pooling layer có thể xác định các vùng genomic quan trọng nhất cho việc phân loại lối sống, giúp sinh học gia khám phá các motifs hoặc genes đặc trưng cho từng lifestyle. Trực quan hóa các feature maps từ CNN layers có thể tiết lộ các patterns cục bộ mà mô hình học được, trong khi gradient-based methods như Integrated Gradients hoặc attention rollout có thể highlight các nucleotides cụ thể đóng góp vào predictions. Những phân tích này không chỉ tăng độ tin cậy của mô hình mà còn có thể cung cấp biological insights mới về các yếu tố quyết định lối sống thực khuẩn thể, potentially leading to discovery of novel genetic markers hoặc regulatory mechanisms.

Cuối cùng, PhaBERT-CNN có tiềm năng được mở rộng sang các tác vụ sinh tin học liên quan khác. Kiến trúc lai kết hợp genomic foundation model với task-specific components có thể được áp dụng cho host prediction—dự đoán vi khuẩn chủ của phage dựa trên genome sequence—bằng cách thay đổi classification head và fine-tune trên tập dữ liệu phage-host pairs. Tương tự, phương pháp có thể được chuyển giao sang prophage detection—xác định các vùng prophage trong bacterial genomes—thông qua sliding window approach và sequence labeling. Việc tích hợp PhaBERT-CNN vào các pipeline phân tích metagenomic tự động như VirSorter hoặc VIBRANT cũng có thể cải thiện độ chính xác tổng thể của các công cụ này. Ngoài ra, tối ưu hóa inference speed thông qua model quantization, knowledge distillation, hoặc efficient attention mechanisms sẽ làm cho PhaBERT-CNN phù hợp hơn cho ứng dụng quy mô lớn trên các tập dữ liệu metagenomic massive. Những hướng phát triển này không chỉ mở rộng khả năng ứng dụng của PhaBERT-CNN mà còn đóng góp vào sự phát triển chung của computational genomics và metagenomic analysis.
